% Generated draft: requires mathematical and source review.
\documentclass[11pt]{article}
\usepackage[margin=25mm]{geometry}
\usepackage{fontspec}
\setmainfont{Latin Modern Roman}
\usepackage{amsmath,amssymb,mathtools,unicode-math}
\setmathfont{Latin Modern Math}
\usepackage{enumitem,needspace,xurl,hyperref}
\setlength{\parindent}{0pt}
\setlength{\parskip}{5pt}
\setlength{\emergencystretch}{3em}
\newcommand{\N}{\mathbb{N}}
\newcommand{\Z}{\mathbb{Z}}
\newcommand{\Q}{\mathbb{Q}}
\newcommand{\R}{\mathbb{R}}
\newcommand{\C}{\mathbb{C}}
\newcommand{\F}{\mathbb{F}}
\newcommand{\A}{\mathbb{A}}
\newcommand{\PP}{\mathbb{P}}
\newcommand{\E}{\mathbb{E}}
\newcommand{\eps}{\varepsilon}
\newcommand{\dd}{\,\mathrm{d}}
\begin{document}
\section*{UnsolvedMath ID: 801172}
\textbf{Source problem code:} catalogue entry
\subsection{Definitions and mathematical statement}
Let \(n\geq 1\), let \(\Gamma=\Z^n\), and call a real matrix \(A\in\operatorname{GL}_n(\R)\) expansive if every eigenvalue of \(A\) has modulus greater than \(1\). Use the Fourier transform convention
\[
\widehat f(\xi)=\int_{\R^n}f(x)e^{-2\pi i x\cdot\xi}\,\dd x
\]
for integrable functions, extended to \(L^2(\R^n)\) by continuity.

For a specified exponent \(N>0\), say that \(\psi\in L^2(\R^n)\) has the required decay of order \(N\) if there exist constants \(C,D,R>0\) such that
\[
|\psi(x)|\leq C|x|^{-N}\quad\text{for almost every }|x|\geq R,
\qquad
|\widehat\psi(\xi)|\leq D|\xi|^{-N}\quad\text{for almost every }|\xi|\geq R.
\]
Here \(|\cdot|\) denotes the Euclidean norm. These almost-everywhere bounds express the supplied decay conditions for representatives of \(L^2\) functions; they impose no additional bounds near the origin.

Under the usual finite-generator interpretation, an orthonormal wavelet family for \((A,\Z^n)\) consists of functions \(\psi_1,\ldots,\psi_r\in L^2(\R^n)\), where \(r\geq1\) is an integer, such that
\[
\left\{|\det A|^{j/2}\psi_\ell(A^j x-k):
 j\in\Z,\ k\in\Z^n,\ 1\leq\ell\leq r\right\}
\]
is an orthonormal basis of \(L^2(\R^n)\).

\textbf{Requested characterization.} Characterize the expansive dilations \(A\) for which such an orthonormal wavelet family exists with every generator well-localized: both the generator and its Fourier transform have polynomial decay of some large order \(N>0\), in the sense above. For finitely many generators, their exponents may be replaced by a common minimum exponent.

\textbf{Unresolved qualifications.} The source says ``some large \(N>0\)'' but supplies no numerical threshold, dependence on dimension, or quantifier specifying how large \(N\) must be. Thus the decay condition above is precise for each fixed \(N\), while the intended range of \(N\) in the requested characterization remains unspecified. The finite-generator interpretation is also explicit here as a convention: the supplied phrase ``possibly with multiple generators'' does not state whether infinitely many generators are allowed. The source describes this question as a variant of Problem 3.1, whose statement is not supplied.
\subsection{Short English statement}
With translation lattice \(\Z^n\), determine which expansive matrices admit an orthonormal wavelet basis whose generators and Fourier transforms decay polynomially at infinity. Multiple generators are allowed. The meaning of a ``large'' decay exponent and the permitted number of generators require clarification.
\subsection{Sources}
\begin{itemize}
\item Local UnsolvedMath registry, permanent ID \(801172\), exact source problem code \texttt{catalogue entry}, in the AMR Open Problem Lists collection. The supplied record attributes the question to Bownik and Rzeszotnik, with proposed year 2019, and labels it \texttt{open}; this label is supplied registry information, not an independently verified status assessment.
\item Bownik and Rzeszotnik, \emph{Open Problems in Wavelet Theory} (2019), source item Problem 11. Supplied source URL: \url{https://arxiv.org/abs/1905.08230}. The supplied provenance records access on 2026-10-16 and extraction by \texttt{arxiv\_latex}; that access information is supplied provenance and was not independently verified for this draft.
\item The supplied statement identifies the question as a variant of Problem 3.1. No statement of that related problem is included in the supplied record.
\end{itemize}
\end{document}
